% Part: computability
% Chapter: recursive-functions
% Section: halting-problem

\documentclass[../../../include/open-logic-section]{subfiles}

\begin{document}

\olfileid{cmp}{rec}{hlt}
\olsection{مشكلة التوقف}

\emph{مشكلة التوقف} بوجه عام هي مشكلة تقرير ما إذا كان حساب دالة عند
دخل~$n$ يتوقف، أي ينتج نتيجة، إذا أعطينا المواصفة~$e$ لهذه الدالة
القابلة للحساب، كبرنامج مثلًا، وأعطينا العدد~$n$. وقد أثبت آلان تورنغ
في نتيجة شهيرة أن هذه المشكلة نفسها لا يمكن حلها بدالة قابلة للحساب؛
أي إن الدالة
\[
h(e, n) =
\begin{cases}
1 & \text{إذا توقف الحساب $e$ عند الدخل $n$}\\
0 & \text{في غير ذلك،}
\end{cases}
\]
ليست قابلة للحساب.

وفي سياق الدوال العودية الجزئية، يمكن أن يؤدي الفهرس~$e$ المعطى في
مبرهنة كليني للصورة السوية دور مواصفة البرنامج. فإذا كانت $f$ دالة
عودية جزئية، كان كل~$e$ تصدق عنده معادلة مبرهنة الصورة السوية فهرسًا
لـ~$f$. وإذا أعطي عدد~$e$، قررت مبرهنة الصورة السوية أن
\[
\cfind{e}(x) \simeq U(\mu s \; T(e, x, s))
\]
دالة عودية جزئية، وأنه لكل دالة عودية جزئية
$f\colon\Nat\to\Nat$ يوجد $e\in\Nat$ بحيث
$\cfind{e}(x)\simeq f(x)$ لكل~$x\in\Nat$. والواقع أنه لا يوجد لكل
$f$ كهذه فهرس واحد فقط، بل عدد لا نهائي من هذه الفهارس~$e$. وتعرف
\emph{دالة التوقف}~$h$ بـ
\[
h(e, x) =
\begin{cases}
1 & \text{إذا كانت $\cfind{e}(x) \fdefined$}\\
0 & \text{في غير ذلك.}
\end{cases}
\]
% تصحيح مصدر (C226-01): كل عدد طبيعي فهرس في التعداد المعتمد، فحُذف الفرع المستحيل.
لاحظ أن $h(e,x)=0$ إذا وفقط إذا كانت $\cfind{e}(x)\fundefined$.

\begin{thm}
\ollabel{thm:halting-problem}
دالة التوقف~$h$ ليست عودية جزئية.
\end{thm}

\begin{proof}
لو كانت $h$ عودية جزئية، لأمكننا تعريف
\[
d(y) =
\begin{cases}
1 & \text{إذا كان $h(y, y) = 0$}\\
\umin{x}{x \neq x} & \text{في غير ذلك.}
\end{cases}
\]
ولما لم يوجد عدد $x$ يحقق $x\neq x$، فلا يوجد
$\umin{x}{x\neq x}$، ومن ثم تكون $d(y)\fundefined$ إذا وفقط إذا كان
$h(y,y)\neq0$. ويترتب على هذا التعريف ما يأتي:
\begin{enumerate}
\item $d(y)\fdefined$ إذا وفقط إذا كان $h(y,y)=0$، أي إذا وفقط إذا
  كانت $\cfind{y}(y)\fundefined$.
\item $d(y)\fundefined$ إذا وفقط إذا كان $h(y,y)=1$، أي إذا وفقط إذا
  كانت $\cfind{y}(y)\fdefined$.
\end{enumerate}
ولو كانت $h$ عودية جزئية، لكانت $d$ عودية جزئية أيضًا. وعندئذ، بحسب
مبرهنة كليني للصورة السوية، يكون لها فهرس~$e_d$، أي
$\cfind{e_d}\simeq d$. لكن التعادلين السابقين يعطيان
\[
d(e_d)\fdefined \quad\text{إذا وفقط إذا}\quad
\cfind{e_d}(e_d)\fundefined
\quad\text{إذا وفقط إذا}\quad d(e_d)\fundefined\text{،}
\]
وهو محال. فلا تكون~$h$ عودية جزئية.
\end{proof}

\end{document}
